\section{USAMO 2009}
        \begin{enumerate}
        	\item (\textit{USAMO 2009}) Given circles $\omega_1$ and $\omega_2$ intersecting at points $X$ and $Y$, let $\ell_1$ be a line through the center of $\omega_1$ intersecting $\omega_2$ at points $P$ and $Q$ and let $\ell_2$ be a line through the center of $\omega_2$ intersecting $\omega_1$ at points $R$ and $S$.  Prove that if $P, Q, R$ and $S$ lie on a circle then the center of this circle lies on line $XY$.


 

	\item (\textit{USAMO 2009}) Let $n$ be a positive integer.  Determine the size of the largest subset of $\{ - n, - n + 1, \ldots , n - 1, n\}$ which does not contain three elements $a, b, c$ (not necessarily distinct) satisfying $a + b + c = 0$.


 

	\item (\textit{USAMO 2009}) We define a \textit{chessboard polygon} to be a polygon whose sides are situated along lines of the form $x = a$ or $y = b$, where $a$ and $b$ are integers. These lines divide the interior into unit squares, which are shaded alternately grey and white so that adjacent squares have different colors. To tile a chessboard polygon by dominoes is to exactly cover the polygon by non-overlapping $1 \times 2$ rectangles. Finally, a \textit{tasteful tiling} is one which avoids the two configurations of dominoes shown on the left below. Two tilings of a $3 \times 4$ rectangle are shown; the first one is tasteful, while the second is not, due to the vertical dominoes in the upper right corner.


 
                    \begin{figure}[H]
                        \centering
                        \includegraphics[width=0.5\textwidth]{Problem/USAMO_I_2009_3_Problem_Figure_1.png}
                    \end{figure}
                    

  a) Prove that if a chessboard polygon can be tiled by dominoes, then it can be done so tastefully.


 b) Prove that such a tasteful tiling is unique.


 

	\item (\textit{USAMO 2009}) For $n \ge 2$ let $a_1$, $a_2$, ..., $a_n$ be positive real numbers such that 


 $(a_1+a_2+ ... +a_n)\left( {1 \over a_1} + {1 \over a_2} + ... +{1 \over a_n} \right) \le \left(n+ {1 \over 2} \right) ^2$

 Prove that $\text{max}\, (a_1, a_2, ... ,a_n) \le  4\, \text{min}\, (a_1, a_2, ... , a_n)$.


 

	\item (\textit{USAMO 2009}) Trapezoid $ABCD$, with $\overline{AB}||\overline{CD}$, is inscribed in circle $\omega$ and point $G$ lies inside triangle $BCD$.  Rays $AG$ and $BG$ meet $\omega$ again at points $P$ and $Q$, respectively.  Let the line through $G$ parallel to $\overline{AB}$ intersect $\overline{BD}$ and $\overline{BC}$ at points $R$ and $S$, respectively.  Prove that quadrilateral $PQRS$ is cyclic if and only if $\overline{BG}$ bisects $\angle CBD$.


 

	\item (\textit{USAMO 2009}) Let $s_1, s_2, s_3, \ldots$ be an infinite, nonconstant sequence of rational numbers, meaning it is not the case that $s_1 = s_2 = s_3 = \cdots.$  Suppose that $t_1, t_2, t_3, \ldots$ is also an infinite, nonconstant sequence of rational numbers with the property that $(s_i - s_j)(t_i - t_j)$ is an integer for all $i$ and $j$.  Prove that there exists a rational number $r$ such that $(s_i - s_j)r$ and $(t_i - t_j)/r$ are integers for all $i$ and $j$.


 

\end{enumerate}